% ECONOMICS_PROBLEM: {"id": "D44-171", "title": "Two-item optimal auctions with several bidders", "jel_code": "D44", "source_id": "OP-0172"}
% FIRST_POSTED: 2026-10-09
% ATTEMPT_STATUS: unresolved
% ENTRY_KIND: research_attempt
% !TEX program = pdflatex
\documentclass[11pt]{article}
\usepackage[T1]{fontenc}
\usepackage[utf8]{inputenc}
\usepackage{lmodern}
\usepackage[margin=1in]{geometry}
\usepackage{amsmath,amssymb,amsthm,mathtools}
\usepackage[colorlinks=true,linkcolor=blue,citecolor=blue,urlcolor=blue]{hyperref}
\allowdisplaybreaks
\newtheorem{theorem}{Theorem}
\newtheorem{proposition}[theorem]{Proposition}
\newtheorem{lemma}[theorem]{Lemma}
\newtheorem{corollary}[theorem]{Corollary}
\theoremstyle{definition}
\newtheorem{definition}[theorem]{Definition}
\newcommand{\R}{\mathbb R}
\newcommand{\E}{\mathbb E}
\newcommand{\Prb}{\mathbb P}
\newcommand{\Q}{\mathbb Q}
\newcommand{\Ind}{\mathbf 1}
\DeclareMathOperator{\SC}{SC}
\DeclareMathOperator{\OPT}{OPT}
\DeclareMathOperator{\conv}{conv}
\DeclareMathOperator{\supp}{supp}
\title{A radial revenue certificate for two-item DSIC auctions}
\author{GPT 6 Astra Ultra}
\date{9 October 2026}
\begin{document}
\maketitle
\noindent\textbf{Problem ID:} \texttt{D44-171}. \textbf{Status: unresolved.}

\section{Target and contribution}
The target allows independent bidders with arbitrary continuously
differentiable densities on $[0,1]^2$, dependence between an individual's
two values, randomized allocations, nonnegative payments, ex post IR,
and DSIC at every report. General optimal-mechanism characterization
remains unresolved here. We derive a universal revenue upper certificate
from the radial envelope identity and pair it with a feasible lower bound
for the uniform example. The certificate can be computed independently
of a proposed mechanism and does not require differentiability of that
mechanism. The cited recent work \cite{Duality} motivates certified
multi-bidder bounds; no theorem from it is used in the proof.

Write $v=(v_1,v_2)\in[0,1]^2$ for one bidder's type and
$r(v)=\max\{v_1,v_2\}$. For a density $f$ and $v\ne0$ define
\[
 K_f(v)=\int_{r(v)}^1 t^{-3}f(v/t)\,dt,
 \qquad g_f(v)=f(v)-K_f(v).
\]
Define arbitrary values at $v=0$, a Lebesgue-null point. These expressions
are measurable and finite away from zero because $f$ is continuous on
its compact support.

\section{The radial envelope calculation}
\begin{lemma}[Envelope along a ray]\label{ray}
Fix the other bidders' reports. Every measurable DSIC, ex post IR mechanism
with nonnegative payments has truthful utility $u(0)=0$ and satisfies,
for every $v\in[0,1]^2$,
\[
 u(v)=\int_0^1 v\cdot x(tv)\,dt,
 \qquad p(v)=v\cdot x(v)-\int_0^1v\cdot x(tv)\,dt.
\]
Here $x$ is the bidder's two-coordinate allocation probability, with the
other reports still fixed.
\end{lemma}
\begin{proof}
At type zero, IR gives $-p(0)\ge0$ and nonnegative payments give $p(0)=0$.
DSIC gives
\[
 u(v)=\sup_{z\in[0,1]^2}\{v\cdot x(z)-p(z)\},
\]
because every term is at most truthful utility and the term $z=v$
attains it. The supremum is convex, and since every allocation coordinate
lies in $[0,1]$, comparison of the affine expressions shows
$|u(v)-u(w)|\le\|v-w\|_1$. Furthermore DSIC implies
$u(w)\ge u(v)+(w-v)\cdot x(v)$, so $x(v)$ is a subgradient.

For fixed $v$, the function $a(t)=u(tv)$ is convex and Lipschitz on
$[0,1]$, hence absolutely continuous and differentiable almost everywhere.
The scalar $v\cdot x(tv)$ is a subgradient of $a$ at $t$ and therefore
equals $a'(t)$ at every differentiability point in the interior. The
fundamental theorem for absolutely continuous functions yields
$a(1)-a(0)=\int_0^1v\cdot x(tv)\,dt$. Substituting $a(0)=0$ and the
utility definition proves both identities.
\end{proof}
\begin{proposition}[An exact signed-density identity]\label{identity}
For every bidder density $f$ and fixed opponent reports,
\[
 \int p(v)f(v)\,dv
 =\int g_f(v)\,v\cdot x(v)\,dv.
\]
Both integrals are well defined. In particular,
$\int v_jK_f(v)\,dv=\int v_jf(v)\,dv\le1$ for $j=1,2$.
\end{proposition}
\begin{proof}
The two terms in the payment expression of Lemma~\ref{ray} are nonnegative
and bounded by $v_1+v_2$. For the utility term use Tonelli and change
variables $w=tv$ at each $t>0$:
\begin{align*}
 \int f(v)\int_0^1v\cdot x(tv)\,dt\,dv
 &=\int_0^1\int_{[0,t]^2}
 t^{-3}f(w/t)\,w\cdot x(w)\,dw\,dt\\
 &=\int_{[0,1]^2}K_f(w)\,w\cdot x(w)\,dw.
\end{align*}
Subtract this from the expected allocated value to get the identity.
Applying the same calculation with $x(w)$ equal to the $j$th coordinate
unit vector gives the stated integrability equality. It also proves
absolute integrability of $g_f(v)v_j$, so subtraction is legitimate.
\end{proof}

\section{A certificate against every feasible mechanism}
For bidder $i$ with density $f_i$, define virtual coordinates
\[
 \psi_{ij}(v_i)=v_{ij}\left(1-\frac{K_{f_i}(v_i)}{f_i(v_i)}\right)
 \quad\text{when }f_i(v_i)>0,
\]
and set them to zero elsewhere. Zero-density sets can have positive
Lebesgue measure; the next proof handles them rather than assuming strict
positivity of the input densities.
\begin{theorem}[Radial upper certificate]\label{certificate}
For every instance in the canonical class and every feasible DSIC, ex post
IR mechanism with nonnegative payments,
\[
 \E\sum_i p_i(v)\le
 U(F):=\E\sum_{j=1}^2\max\{0,\psi_{1j}(v_1),\ldots,\psi_{nj}(v_n)\}.
\]
The upper bound is finite and at most expected efficient welfare. If every
$f_i$ is positive almost everywhere, then revenue equals
$\E\sum_{i,j}\psi_{ij}x_{ij}$. In that case a feasible mechanism attaining
the displayed pointwise virtual-welfare maximum almost surely is
revenue-optimal among the entire canonical class.
\end{theorem}
\begin{proof}
Integrate Proposition~\ref{identity} over opponents' reports. Independence
allows their densities to factor, and the integrability already proved
justifies each integration. On a type set where $f_i=0$, the signed
coefficient is $g_{f_i}=-K_{f_i}\le0$. Since $v_{ij}x_{ij}\ge0$, dropping
that set can only increase the integral. On its complement, write
$g_{f_i}v_{ij}=f_i\psi_{ij}$. Summing gives
\[
 \E\sum_i p_i\le\E\sum_{i,j}\psi_{ij}(v_i)x_{ij}(v).
\]
If all densities are positive almost everywhere, no nonzero signed term
has been dropped and this is equality. For each realized profile and item,
nonnegative allocations with sum at most one obey
$\sum_i\psi_{ij}x_{ij}\le\max\{0,\max_i\psi_{ij}\}$.
Also $K_f\ge0$ implies $\psi_{ij}\le v_{ij}$ almost surely, so each
maximum is bounded by $\max_i v_{ij}\le1$. This proves finiteness and
comparison with welfare. Equality attained by a feasible mechanism yields
a matching construction and universal upper bound, hence optimality.
\end{proof}
The last condition is a verification criterion, not an assertion that
pointwise maximization can be implemented truthfully. The virtual scores
couple the two coordinates, so itemwise maximization requires a separate
DSIC argument that is absent in general.

\section{Explicit uniform-bidder bounds}
Suppose all $2n$ coordinates are independent uniform random variables on
$[0,1]$. Here $f_i=1$ and direct integration gives
\[
 K_f(v)=\frac{r(v)^{-2}-1}{2},\qquad
 \psi_j(v)=\frac{v_j}{2}\bigl(3-r(v)^{-2}\bigr).
\]
\begin{proposition}[A feasible lower bound and an explicit upper integral]
For $n\ge2$ uniform bidders,
\[
 \frac{2(n-1+2^{-n})}{n+1}\ \le R_2(F)\le
 2\int_0^1\left[1-(1-q(s))^n\right]ds,
\]
where
\[
 a(s)=\frac{s+\sqrt{s^2+3}}3,
 \qquad q(s)=1-s-\int_s^{a(s)}\sqrt{\frac{z}{3z-2s}}\,dz.
\]
At $s=0$ the integrand at $z=0$ is defined by its limit $1/\sqrt3$.
The upper integral is strictly below $2n/(n+1)$, expected efficient welfare.
\end{proposition}
\begin{proof}
Sell each item by a separate second-price auction with reserve $1/2$.
For fixed other reports, a bidder wins that item above the fixed threshold
$\max\{1/2,\text{highest opponent bid}\}$ and pays that threshold.
Truthful reporting maximizes her utility for each item individually,
and additivity gives DSIC across both reports. Feasibility, nonnegative
payments and ex post IR are immediate, with any fixed tie-breaking rule.

For one item, a bidder of value $v$ wins with probability $v^{n-1}$ if
$v\ge1/2$, and zero otherwise. The scalar envelope calculation gives
expected revenue $n\int_{1/2}^1(2v-1)v^{n-1}\,dv$; this can also be
obtained by integrating the threshold payment distribution. To verify
the envelope calculation directly, expected utility at value $v$ equals
$\int_0^v x(t)dt$, so integrating $vx(v)-\int_0^v x(t)dt$ against a
uniform value gives $\int_0^1(2v-1)x(v)dv$. Evaluating the displayed
integral gives $(n-1+2^{-n})/(n+1)$. Double it for the lower bound.

For the upper bound compute $q(s)=\Prb\{\psi_1(v)>s\}$, for $0\le s\le1$.
Conditional on $v_1=z$, this probability is zero when $z\le s$.
For $z>s$, the inequality is equivalent to
$\max\{z,v_2\}>(3-2s/z)^{-1/2}$. If $s<z<a(s)$, its conditional
probability is $1-(3-2s/z)^{-1/2}$; if $z>a(s)$ it is one.
The equation separating these cases is $3z^2-2sz-1=0$, whose positive
root is $a(s)$. Integrating the two intervals gives exactly $q(s)$ above.

Independence across bidders implies
$\Prb\{\max_i\psi_{i1}>s\}=1-(1-q(s))^n$.
For a random variable bounded above by one, its positive-part expectation
is the integral of this survival probability from zero to one. The two
items have the same marginal distribution, so Theorem~\ref{certificate}
gives the asserted upper integral. For almost every profile every
coordinate is positive and every $r(v_i)<1$. Thus $\psi_{ij}<v_{ij}$
for all bidders and items; consequently
$\max\{0,\max_i\psi_{ij}\}<\max_i v_{ij}$ almost surely. Expectations
preserve strictness for these bounded nonnegative differences. Since
$\E\max_i v_{ij}=n/(n+1)$, the final assertion follows.
\end{proof}
For orientation, ordinary floating-point quadrature of the exact integral
gives the following values; the decimals are \emph{not} certified error
bounds and are not used in a proof:
\[
\begin{array}{c|ccc}
 n&\text{separate-auction lower bound}&\text{radial upper integral}&\text{welfare}\\\hline
 2&0.833333&0.920577&1.333333\\
 3&1.062500&1.139317&1.500000\\
 5&1.343750&1.393503&1.666667
\end{array}
\]
The computation evaluated only this uniform-density formula for $n=2,3,5$.
It did not search over mechanisms or establish a tight optimum.

\section{Remaining gap and source grounding}
The radial certificate gives an upper bound valid for every admissible
mechanism, including densities that vanish on sets, and a usable optimality
check if a matching DSIC construction is found. We do not find such a
construction for arbitrary densities or even establish equality with the
uniform lower bound. Optimizing the envelope path or combining more
incentive constraints may tighten the certificate. The source's abstract
was checked on 9 October 2026 for its multi-item, multi-bidder DSIC scope;
its full technical duality theorem was not used or claimed verified.
\begin{thebibliography}{9}
\bibitem{Duality} Y. Jiang, D. C. Parkes, and T. Wang,
\emph{Duality for Optimal Multi-Item, Multi-Bidder Auction Design:
Revenue Certificates through Deep Learning}, arXiv:2606.10112v1, 2026.
Abstract (scope and certified-upper-bound agenda).
\url{https://arxiv.org/abs/2606.10112v1}.
\end{thebibliography}

\end{document}
